% ===== Джерела до розділу 5. Ключі з префіксом ch05: =====
% Позиції ядра (bib-core.tex) тут НЕ дублюються.

% --- Формати та внутрішня будова ---
\bibitem{ch05:pecoff} Microsoft. PE Format Specification (Portable Executable and Common Object File Format). Microsoft Docs, 2024. URL: \url{https://learn.microsoft.com/windows/win32/debug/pe-format}.

\bibitem{ch05:elf} Tool Interface Standard (TIS). Executable and Linkable Format (ELF) Specification. Version~1.2. TIS Committee, 1995.

\bibitem{ch05:russinovich} Russinovich M., Solomon D., Ionescu A., Allievi A. Windows Internals. 7th ed. Parts~1--2. Redmond: Microsoft Press, 2017--2021.

\bibitem{ch05:eagle} Eagle C. The IDA Pro Book. 2nd ed. San Francisco: No Starch Press, 2011. 672~p.

% --- Ін'єкції та виконання ---
\bibitem{ch05:barnea} Hosseini A. Ten Process Injection Techniques: A Technical Survey of Common and Trending Process Injection Techniques. Elastic Security Research, 2017.

\bibitem{ch05:klein} Klein T., Kotler I. Process Injection Techniques: Gotta Catch Them All. \emph{Black Hat USA}. Las Vegas, 2019.

\bibitem{ch05:attck-t1055} MITRE ATT\&CK. Technique T1055: Process Injection. The MITRE Corporation. URL: \url{https://attack.mitre.org/techniques/T1055/}.

% --- Персистентність та руткіти ---
\bibitem{ch05:hoglund} Hoglund G., Butler J. Rootkits: Subverting the Windows Kernel. Boston: Addison-Wesley, 2005. 352~p.

\bibitem{ch05:matrosov} Matrosov A., Rodionov E., Bratus S. Rootkits and Bootkits: Reversing Modern Malware and Next Generation Threats. San Francisco: No Starch Press, 2019. 504~p.

\bibitem{ch05:blunden} Blunden B. The Rootkit Arsenal: Escape and Evasion in the Dark Corners of the System. 2nd ed. Burlington: Jones \& Bartlett Learning, 2012. 784~p.

\bibitem{ch05:attck-persist} MITRE ATT\&CK. Tactic TA0003: Persistence. The MITRE Corporation. URL: \url{https://attack.mitre.org/tactics/TA0003/}.

% --- Геші та ідентифікація ---
\bibitem{ch05:imphash} Mandiant. Tracking Malware with Import Hashing (imphash). Mandiant Threat Research, 2014.

% --- Українські джерела ---
\bibitem{ch05:windows-vuln} Дослідження методів виявлення вразливостей програмного забезпечення та ІКС на базі ОС Windows із розробленням засобу автоматизованого аналізу. \emph{Кібербезпека: освіта, наука, техніка}. 2026. URL: \url{https://csecurity.kubg.edu.ua/index.php/journal/article/view/1302}.

\bibitem{ch05:apt-analysis} Аналіз кібератак та діяльності APT груп в Україні. \emph{Кібербезпека: освіта, наука, техніка}. 2024. URL: \url{https://csecurity.kubg.edu.ua/index.php/journal/article/view/574}.
